\documentclass[tikz,border=8pt]{standalone}
\usepackage{tikz}
\usepackage{amsmath}
\usepackage{amssymb}
\usepackage{ctex}
\usetikzlibrary{calc, positioning, arrows.meta, decorations.pathreplacing}

\begin{document}
\begin{tikzpicture}[>=Stealth]

%% ===== 左图：加法封闭性 =====
\begin{scope}[xshift=0cm]

  % 子空间背景（一条过原点的直线）
  \fill[blue!8] (-0.4,-0.6) -- (3.8,5.7) -- (4.2,5.7) -- (0.4,-0.6) -- cycle;
  \draw[blue!40, thick] (-0.4,-0.6) -- (3.8, 5.7)
    node[above right, font=\small, blue!70!black] {子空间 $W$};

  % 坐标轴
  \draw[->, gray!60, thin] (-0.5,0) -- (4.5,0) node[right, font=\scriptsize] {$x$};
  \draw[->, gray!60, thin] (0,-0.5) -- (0, 5.8) node[above, font=\scriptsize] {$y$};

  % 三个向量（都在直线上）
  \draw[->, very thick, blue!80!black] (0,0) -- (1,1.5)
    node[left, font=\small] {$\mathbf{u}$};
  \draw[->, very thick, teal!80!black] (0,0) -- (2,3)
    node[right, font=\small] {$\mathbf{v}$};
  \draw[->, very thick, red!70!black] (0,0) -- (3,4.5)
    node[right, font=\small] {$\mathbf{u}+\mathbf{v}$};

  % 平移箭头（u 从 v 终点出发）
  \draw[->, blue!60, dashed, thick] (2,3) -- (3,4.5);

  % 原点
  \fill (0,0) circle (1.8pt) node[below left, font=\scriptsize] {$O$};

  % 封闭性标注
  \node[font=\small, fill=white, draw=gray!50, thin, rounded corners=2pt, inner sep=3pt]
    at (2.0, -1.0) {\textbf{加法封闭性}};
  \node[font=\scriptsize, fill=white, draw=gray!40, thin, rounded corners=2pt, inner sep=2pt]
    at (2.0, -1.65) {$\mathbf{u},\mathbf{v}\in W \Rightarrow \mathbf{u}+\mathbf{v}\in W$};

\end{scope}

%% ===== 右图：标量乘法封闭性 =====
\begin{scope}[xshift=6.5cm]

  % 子空间背景
  \fill[blue!8] (-0.4,-2.2) -- (2.6,5.7) -- (3.0,5.7) -- (0.4,-2.2) -- cycle;
  \draw[blue!40, thick] (-0.4,-2.2) -- (2.8, 5.7)
    node[above right, font=\small, blue!70!black] {子空间 $W$};

  % 坐标轴
  \draw[->, gray!60, thin] (-0.5,0) -- (4.0,0) node[right, font=\scriptsize] {$x$};
  \draw[->, gray!60, thin] (0,-2.0) -- (0, 5.8) node[above, font=\scriptsize] {$y$};

  % 原向量 v = (1, 2)
  \draw[->, very thick, teal!80!black] (0,0) -- (1,2)
    node[right, font=\small] {$\mathbf{v}$};

  % 0v = 0
  \fill[orange!80!black] (0,0) circle (2.5pt);
  \node[font=\scriptsize, orange!80!black] at (0.5, 0.15) {$0\mathbf{v}=\mathbf{0}$};

  % 1v = v (已画)
  \node[font=\scriptsize, teal!70!black] at (1.55, 1.0) {$1\mathbf{v}$};

  % -1v
  \draw[->, very thick, orange!80!black] (0,0) -- (-1,-2)
    node[left, font=\small] {$-\mathbf{v}$};

  % 2v
  \draw[->, very thick, red!70!black] (0,0) -- (2,4)
    node[right, font=\small] {$2\mathbf{v}$};

  % 原点
  \fill (0,0) circle (1.8pt) node[below right, font=\scriptsize] {$O$};

  % 封闭性标注
  \node[font=\small, fill=white, draw=gray!50, thin, rounded corners=2pt, inner sep=3pt]
    at (1.5, -1.0) {\textbf{标量乘法封闭性}};
  \node[font=\scriptsize, fill=white, draw=gray!40, thin, rounded corners=2pt, inner sep=2pt]
    at (1.5, -1.65) {$c\in\mathbb{R},\,\mathbf{v}\in W \Rightarrow c\mathbf{v}\in W$};

\end{scope}

%% ===== 底部公理框 =====
\node[font=\small, fill=white, draw=gray!50, thin, rounded corners=3pt, inner sep=8pt,
      align=left] at (5.25, -3.6) {
  \textbf{向量空间 8 大公理（$\mathbf{u},\mathbf{v},\mathbf{w}\in V$，$c,d\in\mathbb{F}$）}\\[2pt]
  (A1) $\mathbf{u}+\mathbf{v}=\mathbf{v}+\mathbf{u}$ \quad
  (A2) $(\mathbf{u}+\mathbf{v})+\mathbf{w}=\mathbf{u}+(\mathbf{v}+\mathbf{w})$\\
  (A3) 存在零向量 $\mathbf{0}$ 使 $\mathbf{v}+\mathbf{0}=\mathbf{v}$ \quad
  (A4) 存在 $-\mathbf{v}$ 使 $\mathbf{v}+(-\mathbf{v})=\mathbf{0}$\\
  (S1) $1\cdot\mathbf{v}=\mathbf{v}$ \quad
  (S2) $c(d\mathbf{v})=(cd)\mathbf{v}$\\
  (S3) $c(\mathbf{u}+\mathbf{v})=c\mathbf{u}+c\mathbf{v}$ \quad
  (S4) $(c+d)\mathbf{v}=c\mathbf{v}+d\mathbf{v}$
};

%% 整体标题
\node[font=\large\bfseries] at (5.25, 6.5)
  {向量空间：封闭性公理的几何意义};

\end{tikzpicture}
\end{document}
